% Section 10.5, from lectures/L10/appendix/d_exercises.md.

\section{Övningar}
\label{c10:sec:exercises}\label{c10:app:d}

\begin{aside}
\textbf{Så kontrollerar du ditt arbete.} Övningarna märkta \textbf{Program} skrivs i Microchip
Studio och kontrolleras i simulatorn: förutsäg vad som ska hända, stega, och jämför. Knappar
simuleras genom att biten i \code{PIND} ändras för hand i fönstret I/O (\secref{studio:4-5} i
bilaga~\ref{app:studio}).

Övningarna märkta \textbf{Räkna för hand} och \textbf{Förståelse} görs på papper, utan
miniräknare, och kontrolleras mot bilaga~\ref{app:answers}. Gör varje övning innan du läser
lösningen. Där en övning har ett troligt felsvar tar lösningen upp det också, och säger varför det
är fel.

Instruktionerna finns på referensbladet i bilaga~\ref{app:avr}.

Varje övning är märkt med sin sort. Exakt en övning är en \textbf{Kontroll}, och i det här kapitlet
är det övning~\ref{c10:ex:11}.
\end{aside}

\exercise{Tecken och koder}{Räkna för hand}
\label{c10:ex:1}
Använd tabellen i \secref{c10:a1}.
\begin{parts}
  \item Vilka ASCII-koder har tecknen \code{'5'}, \code{'9'}, \code{'E'} och \code{'F'}? Svara
    hexadecimalt.
  \item Vilket tal blir kvar när ASCII-koden för \code{'8'} minskas med ASCII-koden för
    \code{'0'}?
  \item Förklara med en mening varför omvandlingen av värdena 10--15 behöver sju extra, men inte
    omvandlingen av 0--9.
  \item Raden \code{ldi r16, 'E'} och raden \code{ldi r16, 0x45} assembleras till exakt samma
    maskinkod. Varför skriver man ändå ibland den första?
\end{parts}

\exercise{En byte blir text}{Räkna för hand}
\label{c10:ex:2}
Programmet i \secref{c10:a3} körs med \code{ldi r16, 0x9E} i stället för \code{0x3C}.
\begin{parts}
  \item Fyll i värdet i \code{r24} efter varje instruktion för den höga nibblen: \code{mov},
    \code{swap}, \code{andi}, och sedan det som \code{cpi}/\code{brlo} bestämmer. Vilket värde får
    \code{r20}?
  \item Gör samma sak för den låga nibblen. Vilket värde får \code{r21}?
  \item Vilken text står i minnet om programmet avslutas med \code{sts 0x0100, r20} och
    \code{sts 0x0101, r21}?
  \item Gör om a)--c) för \code{ldi r16, 0x5A}.
\end{parts}

\exercise{Från tecken till värde}{Räkna för hand}
\label{c10:ex:3}
Använd koden i \secref{c10:a4}.
\begin{parts}
  \item Följ tecknet \code{'B'} (\code{0x42}) genom koden. Vilket värde får \code{r22}?
  \item Följ tecknet \code{'4'}.
  \item \emph{(fördjupning)} Någon skriver ett litet \code{b} på tangentbordet, ASCII \code{0x62}.
    Vilket värde lämnar koden i \code{r22}, och varför är det fel? Föreslå en ändring som hanterar
    små bokstäver.
\end{parts}

\exercise{Registerpar}{Räkna för hand}
\label{c10:ex:4}
\begin{parts}
  \item Skriv talen 2500, 256, 255 och 40\,000 som 16-bitars tal i hexadecimal form, och ange vad
    \code{low()} och \code{high()} blir för vart och ett.
  \item \code{r25} innehåller \code{0x12} och \code{r24} innehåller \code{0x34}. Vilket decimalt
    tal är \code{r25:r24}?
  \item Vilket är det största tal ett registerpar kan hålla? Varför just det?
\end{parts}

\exercise{Addition och subtraktion med 16 bitar}{Räkna för hand}
\label{c10:ex:5}
\begin{parts}
  \item \code{r25:r24} = \code{0x12F0} och \code{r23:r22} = \code{0x0A20}. Följ
    \code{add r24, r22} och \code{adc r25, r23} steg för steg, som i \secref{c10:a7}: vad blir
    \code{r24}, \code{C} och \code{r25}? Kontrollera svaret genom att räkna om båda talen och
    summan decimalt.
  \item Någon skriver \code{add r25, r23} i stället för \code{adc r25, r23}. Vilket svar blir det
    då, och med hur mycket är det fel?
  \item \code{r21:r20} = \code{0x0100}, alltså 256. Följ \code{ldi r22, 1}, \code{ldi r23, 0},
    \code{sub r20, r22} och \code{sbc r21, r23}. Vad blir \code{r20}, \code{C} och \code{r21}, och
    vilket tal är resultatet?
  \item Någon vill göra samma sak kortare, med \code{sbiw r20, 1}. Varför går det inte att
    assemblera, och vilket registerpar hade fungerat?
\end{parts}

\exercise{En fördröjning på 10 ms}{Räkna för hand}
\label{c10:ex:6}
\begin{parts}
  \item Hur många klockcykler är 10~ms vid 16~MHz?
  \item En enda 16-bitars loop med \code{sbiw} och \code{brne}, som i \secref{c10:a10}, plus de två
    \code{ldi} före den. Hur många varv \code{N} ska loopen gå för att hela fördröjningen ska bli så
    nära 10~ms som möjligt, och hur många cykler blir det exakt?
  \item Samma fördröjning med subrutinen \code{delay_ms} från \secref{c10:c4}: vilket värde ska
    \code{r24} ha, och hur många cykler tar anropet?
  \item Vilken av de två är närmast 10~ms? Vilken skulle du välja i ett program, och varför?
\end{parts}

\exercise{Knappen som läses baklänges}{Förståelse}
\label{c10:ex:7}
En knapp är kopplad mellan PD2 och jord, som i \secref{c10:b2}.
\begin{parts}
  \item Förklara, med pull-up-motståndet, varför \code{PIND} bit~2 läses som 0 när knappen är
    nedtryckt.
  \item Programmet glömmer att slå på pull-upen (\code{sbi PORTD, 2}). Vad läser programmet när
    knappen är släppt? Varför är det svårt att hitta det felet genom att prova?
  \item Varför innehåller \code{buttons_to_leds.asm} instruktionen \code{com}? Vad hade lysdioderna
    visat utan den, när ingen knapp är nedtryckt?
\end{parts}

\exercise{Följ programmet}{Räkna för hand}
\label{c10:ex:8}
Programmet i \secref{c10:b3} läser \code{PIND} och visar det inverterade värdet på \code{PORTB}.
\begin{parts}
  \item \code{PIND} läses som \code{0b11111011}. Vilka knappar är nedtryckta, och vad blir
    \code{PORTB}?
  \item Programmet ändras så att \code{andi r16, 0xF0} läggs in mellan \code{com} och \code{out}.
    Vad blir \code{PORTB} nu för samma \code{PIND}? Och för \code{PIND} = \code{0b01101111}?
  \item Hur många gånger per sekund ungefär läser loopen porten vid 16~MHz? Räkna cyklerna för de
    fyra instruktionerna i loopen.
\end{parts}

\exercise{Start och stopp med självhållning}{Program}
\label{c10:ex:9}
I kapitel~\ref{c2:ch} byggdes en start/stopp-krets med ett relä som håller sig självt: ett tryck på
startknappen drar reläet, och det förblir draget när knappen släpps, tills stoppknappen trycks.
Samma funktion ska nu skrivas som ett program.
\begin{itemize}
  \item Startknappen sitter på PD2 och stoppknappen på PD3, båda mot jord med pull-up.
  \item En \enquote{motor}, en lysdiod på PB0, startar när startknappen trycks, och fortsätter gå
    när knappen släpps.
  \item Motorn stannar när stoppknappen trycks.
  \item Om båda knapparna hålls nedtryckta samtidigt ska motorn \textbf{stå still}: stopp går före
    start, av säkerhetsskäl.
\end{itemize}
\begin{parts}
  \item Rita en flödesplan innan du skriver någon kod.
  \item Skriv programmet, med \code{sbic} eller \code{sbis}.
  \item Testa i simulatorn alla fyra fallen: start, släpp start, stopp, och båda samtidigt.
  \item Var i programmet finns \enquote{minnet}, det som motsvarar reläets självhållning?
\end{parts}

\exercise{\code{rcall}, \code{ret} och stacken}{Förståelse}
\label{c10:ex:10}
\begin{parts}
  \item Beskriv de två saker \code{rcall} gör, i den ordning den gör dem.
  \item Varför behöver \code{ret} stacken? Kunde returadressen ha sparats i ett av de 32 registren i
    stället?
  \item \code{SP} är \code{0x08FF} när ett \code{rcall} körs. Vad är \code{SP} när subrutinen
    börjar? Och efter \code{ret}?
  \item Efter \code{ret} ligger returadressen kvar i minnet. Varför är det inget problem?
\end{parts}

\exercise{Kontroll: förutsäg stacken}{Kontroll}
\label{c10:ex:11}
\emph{Förutsäg för hand, kontrollera i simulatorn, förklara skillnaden.}

Programmet nedan anropar en subrutin som lägger två register på stacken. Adresserna i
programminnet kommer från listfilen och står i kommentarerna.

\begin{asmcode}
.include "m328Pdef.inc"

.org 0x0000
    rjmp main                       ; 0x0000

main:
    ldi r16, high(RAMEND)           ; 0x0001
    out SPH, r16                    ; 0x0002
    ldi r16, low(RAMEND)            ; 0x0003
    out SPL, r16                    ; 0x0004
    ldi r16, 0xA5                   ; 0x0005
    ldi r17, 0x3C                   ; 0x0006
    rcall save_two                  ; 0x0007
end:
    rjmp end                        ; 0x0008

save_two:
    push r16                        ; 0x0009
    push r17                        ; 0x000A
here:
    pop r17                         ; 0x000B
    pop r16
    ret
\end{asmcode}

\task{a) För hand}
När programmet står på etiketten \code{here}, alltså efter de två \code{push}: vilket värde har
\code{SP}, och vilka byte ligger på adresserna \code{0x08FF}, \code{0x08FE}, \code{0x08FD} och
\code{0x08FC}? Skriv upp svaret innan du går vidare.

\task{b) I simulatorn}
Skriv in programmet, sätt en brytpunkt på \code{pop r17} och kör dit. Läs \code{SP} i Processor
Status och bytena i fönstret Memory med \textbf{data IRAM} valt.

\task{c) Jämför}
Stämmer din förutsägelse? Den vanligaste avvikelsen är att returadressens två byte står i fel
ordning, eller att \code{SP} är en byte fel. Om ditt svar skiljer sig: vilket av de två felen var
det, och varför?

\task{d)}
Stega vidare till instruktionen efter \code{ret}. Vilket värde har \code{SP} nu, och vilka av de
fyra bytena har ändrats?

\exercise{Hitta felet}{Förståelse, Fördjupning}
\label{c10:ex:12}
\begin{parts}
  \item En subrutin börjar med \code{push r16} och \code{push r17}, och slutar med \code{pop r16},
    \code{pop r17} och \code{ret}. Den assembleras utan fel och återvänder rätt. Vad är fel, och hur
    märks det i programmet som anropar den?
  \item En annan subrutin börjar med \code{push r16} och \code{push r17}, men slutar med bara
    \code{pop r17} och \code{ret}. Vad händer när \code{ret} körs? Varför märks felet mycket
    tydligare än i a)?
  \item Hur kan du i simulatorn, med ett enda värde, se att en subrutin inte lämnar stacken som den
    fann den?
\end{parts}

\exercise{Hela byten som text}{Program, Fördjupning}
\label{c10:ex:13}
Skriv en subrutin \code{byte_to_ascii} som gör om byten i \code{r24} till två tecken: det för den
höga nibblen i \code{r24} och det för den låga i \code{r25}. Den ska använda \code{nibble_to_ascii}
från \secref{c10:c9} och följa kursens regel om register.
\begin{parts}
  \item Vilka register ändrar \code{byte_to_ascii}, och vilka av dem måste sparas?
  \item Skriv subrutinen och ett huvudprogram som anropar den med \code{0xB7} och lägger tecknen på
    adress \code{0x0100} och \code{0x0101}.
  \item Kontrollera i fönstret Memory att texten \code{B7} står där. Hur djup blir stacken som mest?
\end{parts}
